% SEGMENT OLP-0015-S01
% Part: sets-functions-relations
% Chapter: relations
% Section: equivalence-relations

\documentclass[../../../include/open-logic-section]{subfiles}

\begin{document}

\olfileid{sfr}{rel}{eqv}

\olsection{சமானத் தொடர்புகள்}

ஒரு கணத்தின் மீதான ஒன்றாதல் தொடர்பு தற்சுட்டு, சமச்சீர், கடப்பு
ஆகிய பண்புகளைக் கொண்டது. இந்த மூன்று பண்புகளையும் கொண்ட
$R$ என்னும் தொடர்புகள் மிகவும் பொதுவானவை.

% SEGMENT OLP-0015-S02
\begin{defn}[சமானத் தொடர்பு]
தற்சுட்டு, சமச்சீர், கடப்பு ஆகிய பண்புகளைக் கொண்ட
$R \subseteq A^2$ என்னும் தொடர்பு ஒரு \emph{சமானத் தொடர்பு} எனப்படும்.
$A$ இன் !!^{element}s ஆன $x$, $y$ ஆகியவற்றுக்கு $Rxy$
என்பது உண்மையெனில், அவை \emph{$R$-சமானமானவை} எனப்படும்.
\end{defn}

% SEGMENT OLP-0015-S03
சமானத் தொடர்புகளிலிருந்து \emph{சமான வகுப்பு} என்ற கருத்து எழுகிறது.
ஒரு சமானத் தொடர்பு சார்பகத்தை வெவ்வேறு பிரிவுகளாகப் ``பகுக்கிறது''.
ஒவ்வொரு பிரிவிலும் உள்ள அனைத்துப் பொருள்களும் ஒன்றுக்கொன்று தொடர்புடையவை;
வெவ்வேறு பிரிவுகளில் உள்ள எந்தப் பொருள்களும் ஒன்றுக்கொன்று தொடர்புடையவை அல்ல.
சில சமயங்களில் இந்தப் பிரிவுகளைப் பற்றி \emph{நேரடியாகப்} பேசுவது பயனுள்ளதாகும்.
அதற்காக ஒரு வரையறையை அறிமுகப்படுத்துகிறோம்:

% SEGMENT OLP-0015-S04
\begin{defn}\ollabel{def:equivalenceclass}
$R \subseteq A^2$ ஒரு சமானத் தொடர்பாக இருக்கட்டும்.
ஒவ்வொரு $x \in A$ க்கும், $A$ இல் $x$ இன் \emph{சமான வகுப்பு} என்பது
$\equivrep{x}{R} = \Setabs{y \in A}{Rxy}$ என்னும் கணமாகும்.
$R$ ஐப் பொறுத்து $A$ இன் \emph{ஈவுக் கணம்} என்பது
$\equivclass{A}{R} = \Setabs{\equivrep{x}{R}}{x \in A}$ ஆகும்;
அதாவது, இந்தச் சமான வகுப்புகளைக் கொண்ட கணமே அது.
\end{defn}

% SEGMENT OLP-0015-S05
அடுத்த முடிவு, சமான வகுப்புகள் உண்மையில் $A$ இன் பிரிவுகளே என்பதை
நிறுவுவதன் மூலம், சமான வகுப்பின் வரையறையை நியாயப்படுத்துகிறது:

\begin{prop}
$R \subseteq A^2$ ஒரு சமானத் தொடர்பு எனில், $Rxy$ என்பது
$\equivrep{x}{R} = \equivrep{y}{R}$ ஆக இருக்கும்போதும்
அப்போது மட்டுமே உண்மை.
\end{prop}

% SEGMENT OLP-0015-S06
\begin{proof}
இடமிருந்து வலமாகச் செல்லும் திசைக்கு, $Rxy$ எனக் கொண்டு,
$z \in \equivrep{x}{R}$ என்க. அப்போது வரையறையின்படி $Rxz$.
$R$ ஒரு சமானத் தொடர்பு என்பதால் $Ryz$.
(விரிவாகச் சொன்னால்: $Rxy$ ஆகவும் $R$ சமச்சீராகவும் இருப்பதால் $Ryx$;
$Rxz$ ஆகவும் $R$ கடப்புப் பண்புடையதாகவும் இருப்பதால் $Ryz$.)
எனவே $z \in \equivrep{y}{R}$. இதனைப் பொதுமைப்படுத்தினால்,
$\equivrep{x}{R} \subseteq \equivrep{y}{R}$.
ஆனால் இதே முறையில் $\equivrep{y}{R} \subseteq \equivrep{x}{R}$.
எனவே உறுப்புசார் சமத்துவத்தின்படி
$\equivrep{x}{R} = \equivrep{y}{R}$.

% SEGMENT OLP-0015-S07
வலமிருந்து இடமாகச் செல்லும் திசைக்கு,
$\equivrep{x}{R} = \equivrep{y}{R}$ எனக் கொள்வோம்.
$R$ தற்சுட்டுப் பண்புடையது என்பதால் $Ryy$;
ஆகவே $y \in \equivrep{y}{R}$.
$\equivrep{x}{R} = \equivrep{y}{R}$ என்ற எடுகோளின்படி,
$y \in \equivrep{x}{R}$ என்பதும் உண்மை. எனவே $Rxy$.
\end{proof}

% SEGMENT OLP-0015-S08
\begin{ex}
மட்டு எண்கணிதத்திலிருந்து சமானத் தொடர்புகளுக்கு ஓர் அழகிய எடுத்துக்காட்டு கிடைக்கிறது.
எந்த $a$, $b$, $n \in \PosInt$ ஆகியவற்றுக்கும்,
$a$ ஐ $n$ ஆல் வகுக்கும்போது கிடைக்கும் மீதியும்,
$b$ ஐ $n$ ஆல் வகுக்கும்போது கிடைக்கும் மீதியும் ஒன்றாக இருக்கும்போது,
அப்போதும் அப்போது மட்டுமே $a \equiv_n b$ எனக் கூறுவோம்.
(சற்று அதிகக் குறியீடுகளுடன்: ஏதாவது ஒரு $k \in \Int$ க்கு
$a - b = kn$ ஆக இருக்கும்போது, அப்போதும் அப்போது மட்டுமே $a \equiv_n b$.)
இப்போது எந்த $n$ க்கும் $\equiv_n$ ஒரு சமானத் தொடர்பாகும்.
$\equiv_n$ உருவாக்கும் வெவ்வேறு சமான வகுப்புகள் சரியாக $n$ உள்ளன;
அதாவது, $\equivclass{\Nat}{\equiv_n}$ இல் $n$ !!{element}s உள்ளன.
அவை: $n$ ஆல் மீதியின்றி வகுபடும் எண்களின் கணமான
$\equivrep{0}{\equiv_n}$; $n$ ஆல் வகுக்கும்போது மீதி $1$ தரும் எண்களின்
கணமான $\equivrep{1}{\equiv_n}$; \ldots;
$n$ ஆல் வகுக்கும்போது மீதி $n-1$ தரும் எண்களின் கணமான
$\equivrep{n-1}{\equiv_n}$.
\end{ex}

% SEGMENT OLP-0015-S09
\begin{prob}
எந்த $n \in \PosInt$ க்கும் $\equiv_n$ ஒரு சமானத் தொடர்பாகும் என்பதையும்,
$\equivclass{\Nat}{\equiv_n}$ இல் சரியாக $n$ உறுப்புகள் உள்ளன என்பதையும் நிறுவுக.
\end{prob}

\end{document}
